%% =================================================================
%% Chen et al. Measurement 264 (2026) 120300.
%% Reconstruction of page 12 (printed p. 12).
%% Contents: Fig. 16 (FPCB fatigue test, 300-cycle compression depth
%% + microscope before/after), prose tail of Section 4.5 and full
%% Section 4.6 (FPCB defect detection application).
%% =================================================================

\documentclass[11pt]{article}

\usepackage[margin=0.85in]{geometry}
\usepackage{amsmath, amssymb}
\usepackage{mathrsfs}
\usepackage{xcolor}
\usepackage{fancyhdr}

\pagestyle{fancy}
\fancyhf{}
\lhead{\small Z.~Chen et al.}
\rhead{\small Measurement 264 (2026) 120300}
\cfoot{\small 12 $\to$ \arabic{page}}
\renewcommand{\headrulewidth}{0.4pt}

\setcounter{page}{1}
\setcounter{equation}{1}
\setlength{\parskip}{6pt}
\setlength{\parindent}{0pt}

\begin{document}

% ---------- Fig. 16 placeholder ----------
\begin{center}
\fbox{\begin{minipage}[c][8cm][c]{0.95\textwidth}
\centering
\textbf{Fig. 16.}\ Fatigue testing of encapsulated flexible printed
circuit board.\\[4pt]
\textcolor{gray}{\footnotesize
\textbf{(a)} Compression depth of the probe tip at the test point
of the encapsulated FPCB under 300 cycles of scanning: depth vs.
cycle number 0--300, applied pressure 4 MPa, line held at
$\approx 0.18$ mm across all cycles (negligible drift).\\[3pt]
\textbf{(b)} Microscope image of the test point of the
encapsulated FPCB before scanning: intact orange serpentine trace
with a smooth encapsulated surface.\\[3pt]
\textbf{(c)} Microscope image after 300 cycles of scanning: an
almost indiscernible circular indentation mark on the AB-glue
encapsulation layer, no crack or puncture.\\[3pt]
\textbf{(d)} Microscope image of the FPCB itself after 300 cycles
of scanning: the underlying circuit remains intact with no visible
deformation.]}
\end{minipage}}
\end{center}

\vspace{0.4em}

\textbf{Fig. 16.}\ Fatigue testing of encapsulated flexible printed
circuit board. (a) Compression depth of the probe tip at the test
point of the encapsulated FPCB under 300 cycles of scanning.
(b) Microscope image of the test point of the encapsulated FPCB
before scanning. (c) Microscope image of the test point of the
encapsulated FPCB after 300 cycles scanning. (d) Microscope image
of the test point of the FPCB after 300 cycles scanning.

\vspace{0.8em}

% ---------- Prose ----------
scanned by the tactile tomography system with a scanning speed and
scanning step of 2.4 mm/s \& 0.5 mm under a constant pressure of 4
MPa. Fig.~14d, e, and f presents the reconstructed 3D images
corresponding to the soft surface layer of AB glue, PDMS, and
Ecoflex, respectively. All these 3D images successfully
reconstructed the topographical profile of the spherical cap model
beneath the soft surface layer of AB glue, PDMS, and Ecoflex,
indicating that the tactile tomography system has excellent
universality for different materials.

\subsection*{4.6.\ Application in nondestructive testing of encapsulated flexible printed circuit boards (FPCBs)}

Defect Detection of encapsulated flexible printed circuit boards
(FPCBs) is one of the most promising applications for the tactile
tomography system. FPCBs as a popular substrate has been widely
used in flexible wearable devices [39]. In practice, flexible
wearable devices are usually encapsulated with a soft layer to
protect them from the outside environment without sacrificing
their flexibility [40]. However, the current research mainly
focuses on the defect detection of FPCBs without packaging.
Therefore, nondestructive testing of encapsulated FPCBs is still a
challenge. As shown in Fig.~15a, an FPCB sample consisting of PI
and copper lines was purchased, and it was then encapsulated with
a soft layer of AB glue, which was dyed black to prevent detection
by the human eye and optical microscope (Fig.~15b). The
encapsulated FPCBs sample without any defect was scanned by the
tactile tomography system with a scanning speed and scanning step
of 0.6 mm/s \& 0.5 mm under an effective scanning range of $28
\times 10$ mm$^2$ (Fig.~15c). The reconstructed 3D image is shown
in Fig.~15e, in which the layout of the copper lines was
completely reproduced, demonstrating that the tactile tomography
system can detect and image the internal structure of flexible
wearable devices without peeling off the encapsulated layer. To
investigate the defect detection of the encapsulated FPCB, three
defects of open the copper lines on the FPCB were made (as shown
in Fig.~15d). These defects can be found easily on the
reconstructed 3D image in Fig.~15f, meaning that the tactile
tomography system possesses the capability for defect detection in
encapsulated FPCBs. For comparison, ultrasonic imaging (Ultrasonic
detection system, Feli 650) was also employed to inspect the same
encapsulated FPCB sample, the corresponding ultrasonic C-scan
images are presented in Fig.~15g and h. In Fig.~15g, the
ultrasonic image of the intact encapsulated FPCB shows a
relatively blurred circuit pattern, and in Fig.~15h, when defects
exist, the ultrasonic image fails to clearly delineate the defects.
This is because ultrasonic imaging is susceptible to the influence
of the encapsulation layer. The acoustic properties of the
encapsulation material differ from those of the FPCB, causing
scattering and attenuation of ultrasonic waves, which degrades the
imaging quality and makes it difficult to accurately identify
internal defects. This limitation arises from the acoustic
impedance mismatch and the wave-damping nature of the soft

\end{document}
